\documentclass[11pt]{article}

\usepackage[T1]{fontenc}
\usepackage[utf8]{inputenc}
\usepackage{amsmath,amssymb,amsthm,mathtools}
\usepackage{booktabs}
\usepackage{enumitem}
\usepackage{microtype}
\usepackage[hidelinks]{hyperref}
\usepackage[margin=1in]{geometry}

\newtheorem{theorem}{Theorem}[section]
\newtheorem{lemma}[theorem]{Lemma}
\newtheorem{proposition}[theorem]{Proposition}
\newtheorem{corollary}[theorem]{Corollary}
\theoremstyle{definition}
\newtheorem{definition}[theorem]{Definition}
\newtheorem{remark}[theorem]{Remark}

\newcommand{\N}{\mathbb{N}}
\newcommand{\Tsys}{\mathcal{T}}
\newcommand{\Xsys}{\mathcal{X}}
\newcommand{\gec}{\mathrel{\geq_{\mathrm c}}}
\newcommand{\gtc}{\mathrel{>_{\mathrm c}}}
\newcommand{\mat}[4]{\begin{pmatrix}#1&#2\\#3&#4\end{pmatrix}}

\title{Obstructions to Low-Dimensional Natural Matrix Interpretations\\
for a Mixed-Base Collatz Rewrite System}
\author{
Kadyrbekov Khamit\thanks{Corresponding author: \texttt{kadirbekov@gmail.com}}\\
\small Independent Researcher, Tashkent, Uzbekistan\\
\small \texttt{kadirbekov@gmail.com}
\and
Kadirbekov Daniyal\\
\small Independent Researcher, Tashkent, Uzbekistan\\
\small \texttt{daniyal.kadirbekov@gmail.com}\\
\small ORCID: \href{https://orcid.org/0009-0006-0320-9821}{0009-0006-0320-9821}}
\date{}

\begin{document}
\maketitle

\begin{abstract}
Yolcu, Aaronson, and Heule encoded the Collatz conjecture as termination of
an eleven-rule mixed binary--ternary string-rewriting system and asked
whether standard termination interpretations can orient that system.  We
prove a complete negative result for one natural proof language.  No
one-stage relative affine matrix interpretation over the nonnegative
integers exists in dimensions one or two when the two rules implementing
the Collatz update are required to be strict and the nine base-conversion
rules are weak.  In dimension two the matrix entries and affine offsets are
unbounded.  The proof combines exact Farkas certificates, a rigidity lemma
for the six carry rules, a Perron--Frobenius argument for the irreducible
case, and an exhaustive symbolic classification of the remaining zero
patterns.  A deterministic checker reconstructs every parameterized
certificate using integer arithmetic; bounded enumeration was used only to
discover the certificate families and to regression-test the symbolic
proof.  The result does not resolve the Collatz conjecture and does not
exclude higher-dimensional, arctic, tropical, dependency-pair, or
multi-stage interpretations.  It instead settles the unbounded
two-dimensional natural one-stage case of a broader obstruction question
raised by the original encoding, and illustrates a reproducible route from
finite counterexample-guided search to a universal no-go theorem.
\end{abstract}

\noindent\textbf{Keywords:} Collatz conjecture; string rewriting;
termination; matrix interpretation; Farkas certificate; computer-assisted
mathematics

\section{Introduction}

The Collatz map sends an integer $n>0$ to $n/2$ when $n$ is even and to
$3n+1$ when $n$ is odd.  Its conjectured eventual arrival at $1$ remains
open.  Yolcu, Aaronson, and Heule (YAH) translated the conjecture into the
termination of a finite string-rewriting system and applied automatic
termination tools to weakened variants~\cite{yah}.  Their construction is
especially useful because it turns a familiar number-theoretic question
into a sharply specified question about proof search.

Natural matrix interpretations are a standard termination technique
\cite{endmat,zantema}.  A successful interpretation for the full YAH system
would prove its termination and hence Collatz.  Failure in a prescribed
class has the opposite meaning: it rules out a proof language, not the
conjecture.  YAH explicitly raised the possibility that their full system
might admit no suitable matrix or arctic interpretation.  This paper
settles one precise part of that question.

Our principal result is a universal no-go theorem for one-stage relative
natural affine interpretations in dimensions $1$ and $2$.  Only the two
rules that perform the Collatz step are required to decrease strictly; the
nine rules that convert between bases are allowed to be weak.  Thus the
result addresses the rule-removal obligation that actually matters, rather
than the unnecessarily strong requirement that every rule be strict.

The dimension-two proof was found in three stages.  First, exact exhaustive
search over small matrices produced integer dual certificates.  Second,
the supports of these certificates were generalized to parameterized
families.  Finally, the uncovered symbolic region was classified without
any bound on the entries.  Computation is therefore a discovery and checking
device, not the source of the universal quantifier.  The mathematical proof
is given below; the artifact independently expands and checks the
certificate algebra.

To the best of our knowledge, no earlier publication proves the unbounded
dimension-two no-go theorem stated here.  The scope is deliberately narrow:
dimensions at least three, other ordered semirings, dependency pairs, and
sequential rule-removal strategies remain open.

\paragraph{Contributions.}
\begin{enumerate}[leftmargin=*]
  \item A scalar no-go theorem for the relative YAH obligation.
  \item A complete unbounded two-dimensional no-go theorem over $\N$.
  \item A finite collection of parameterized Farkas certificates and a
        deterministic exact checker.
  \item A case study in converting bounded CEGAR-style proof search into a
        symbolic theorem while retaining failed conjectures and scope limits.
\end{enumerate}

\section{The rewrite system and interpretations}

Let the alphabet be
$\Sigma=\{f,t,0,1,2,<,>\}$.  The mixed-base system $\Tsys$ consists of
the rules in Table~\ref{tab:rules}.  We write words by juxtaposition.

\begin{table}[ht]
\centering
\caption{The YAH mixed binary--ternary system $\Tsys$.}
\label{tab:rules}
\begin{tabular}{lll}
\toprule
class & name & rule\\
\midrule
dynamic & even & $f>\to >$\\
dynamic & odd  & $t>\to 2>$\\
\midrule
carry & f0 & $f0\to0f$\\
carry & f1 & $f1\to0t$\\
carry & f2 & $f2\to1f$\\
carry & t0 & $t0\to1t$\\
carry & t1 & $t1\to2f$\\
carry & t2 & $t2\to2t$\\
\midrule
left & 0 & $<0\to<t$\\
left & 1 & $<1\to<ff$\\
left & 2 & $<2\to<ft$\\
\bottomrule
\end{tabular}
\end{table}

YAH prove that termination of $\Tsys$ is equivalent to the Collatz
conjecture~\cite[Theorem~3.17]{yah}.  Let $\Xsys$ denote the nine
non-dynamic rules.  They terminate independently, so a rule-removal proof
may orient the two dynamic rules strictly and every rule in $\Xsys$ weakly.

\begin{definition}[Natural affine matrix interpretation]
Fix $d\ge1$.  Each symbol $s\in\Sigma$ is assigned
\[
  [s](x)=M_sx+v_s,
  \qquad M_s\in\N^{d\times d},\quad v_s\in\N^d.
\]
The upper-left entry of every $M_s$ is positive, the standard sufficient
condition for strict monotonicity in the first component.  The
interpretation of a word is functional composition from left to right.
We use the orders
\[
 x\gec y\iff x_i\ge y_i\ \text{for all }i,
 \qquad
 x\gtc y\iff x_0>y_0\ \text{and }x_i\ge y_i\ (i>0).
\]
An interpretation is \emph{relative} for $(\Tsys,\Xsys)$ if the two dynamic
rules are oriented by $\gtc$ and the other nine by $\gec$.
\end{definition}

For a word $w=s_1\cdots s_k$ write
\[
 M_w=M_{s_1}\cdots M_{s_k},\qquad
 v_w=\sum_{j=1}^{k}M_{s_1}\cdots M_{s_{j-1}}v_{s_j}.
\]
Orientation of $\ell\to r$ for all $x\in\N^d$ requires
$M_\ell\ge M_r$ entrywise and the corresponding weak or strict inequality
between $v_\ell$ and $v_r$.

For dimension two we abbreviate
\[
 F=M_f,\quad T=M_t,\quad Z=M_0,\quad O=M_1,\quad D=M_2,
 \quad L=M_<,\quad G=M_>,
\]
and let $\ell=e_0^{\mathsf T}L=(a,b)$, where $a\ge1$ and $b\ge0$.
The eleven coefficient obligations are
\begin{align}
 FG&\ge G, & TG&\ge DG,\label{eq:dyncoeff}\\
 FZ&\ge ZF, & FO&\ge ZT, & FD&\ge OF,\label{eq:fcarry}\\
 TZ&\ge OT, & TO&\ge DF, & TD&\ge DT,\label{eq:tcarry}\\
 LZ&\ge LT, & LO&\ge LF^2, & LD&\ge LFT.\label{eq:leftcoeff}
\end{align}
All matrix inequalities in this paper are entrywise unless stated otherwise.

\section{Exact dual certificates}

Once the matrices are fixed, the offset obligations are linear inequalities
in the $7d$ nonnegative offset variables.  We use the following elementary
form of Farkas' lemma~\cite{schrijver}.

\begin{lemma}[Certificate criterion]\label{lem:farkas}
For each rule $r$ and component $i$, write the offset difference as
$\Delta_r^{(i)}(v)$ and its required right-hand side as
$\epsilon_r^{(i)}$, where $\epsilon=1$ exactly for component $0$ of a
strict rule and $0$ otherwise.  If nonnegative integers
$\lambda_{r,i}$ satisfy
\[
  \sum_{r,i}\lambda_{r,i}\Delta_r^{(i)}(v)\le0
  \quad\text{for every }v\in\N^{7d},
  \qquad
  \sum_{r,i}\lambda_{r,i}\epsilon_r^{(i)}>0,
\]
then the offset system is infeasible.
\end{lemma}

\begin{proof}
The assumed orientation inequalities imply that the displayed left side is
at least the positive displayed right side.  The coefficient condition and
$v\ge0$ imply that it is at most zero, a contradiction.
\end{proof}

The artifact represents a certificate as a finite list
$(r,i,\lambda_{r,i})$, reconstructs every coefficient from the original
rules, and accepts only if all coefficients are nonpositive and the strict
right side is positive.  No floating-point result is trusted.

\section{The scalar obstruction}

\begin{theorem}\label{thm:scalar}
There is no one-dimensional natural affine interpretation in which the two
dynamic rules of $\Tsys$ are strict and the nine rules of $\Xsys$ are weak.
\end{theorem}

\begin{proof}
Write $[s](x)=a_sx+b_s$ with $a_s\ge1$ and $b_s\ge0$.  The coefficient
parts of five weak obligations give
\[
 a_2\ge a_fa_t,\quad a_t\ge a_2,\quad
 a_fa_1\ge a_0a_t,\quad a_fa_2\ge a_1a_f,\quad
 a_ta_0\ge a_1a_t.
\]
The first two force $a_f=1$ and $a_t=a_2$.  Cancelling positive factors in
the last three gives
$a_1\ge a_0a_2$, $a_2\ge a_1$, and $a_0\ge a_1$; hence
$a_1\ge a_1^2$ and
\[
 a_f=a_t=a_0=a_1=a_2=1.
\]
Strictness of $t>\to2>$ requires $b_t>b_2$, whereas weakness of
$<2\to<ft$ gives $b_2\ge b_f+b_t\ge b_t$.  This is impossible.
\end{proof}

\section{The two-dimensional obstruction}

We now prove the main result.  There is no upper bound on a matrix entry or
offset in this section.

\subsection{The first certificate and its complement}

\begin{lemma}[First-row growth]\label{lem:first}
If \eqref{eq:dyncoeff}--\eqref{eq:leftcoeff} hold and
$\ell F\ge\ell$, then the relative offset system is infeasible.
\end{lemma}

\begin{proof}
The last inequality of \eqref{eq:leftcoeff} gives
$\ell D\ge\ell FT\ge\ell T$.  Take component $0$ of the odd dynamic
rule with weight $a$, component $1$ with weight $b$, and component $0$ of
$<2\to<ft$ with weight $1$.  The combined coefficients of the offsets
$2,f,t,>$ are respectively $0,-\ell,\ell-\ell F,$ and
$\ell T-\ell D$; every other coefficient is zero.  They are all
nonpositive, while the strict right side is $a>0$.  Apply
Lemma~\ref{lem:farkas}.
\end{proof}

\begin{lemma}[Shape of the complement]\label{lem:complement}
If the hypotheses of Lemma~\ref{lem:first} hold but $\ell F\not\ge\ell$,
then
\[
 b>0,\qquad f_{11}=0,\qquad af_{01}<b.
\]
\end{lemma}

\begin{proof}
The first component of $\ell F$ is at least $a$ because $f_{00}\ge1$.
Thus failure occurs in the second component:
$af_{01}+bf_{11}<b$.  This implies all three assertions.
\end{proof}

\subsection{Rigidity of carry inequalities}

\begin{lemma}[Diagonal and determinant rigidity]\label{lem:rigidity}
In each of the six carry inequalities
\eqref{eq:fcarry}--\eqref{eq:tcarry}, the compared products have equal
diagonal entries and equal determinants.  Consequently, if both
off-diagonal entries of the smaller product are positive, the entire
inequality is an equality.
\end{lemma}

\begin{proof}
For $FZ\ge ZF$ and $TD\ge DT$, equality of traces forces equality of the
two diagonal entries.  For the other four rules, cyclic invariance of trace
gives the closed chain
\[
 \operatorname{tr}(FO)\ge\operatorname{tr}(ZT)
 =\operatorname{tr}(TZ)\ge\operatorname{tr}(OT)
 =\operatorname{tr}(TO)\ge\operatorname{tr}(DF)
 =\operatorname{tr}(FD)\ge\operatorname{tr}(OF)
 =\operatorname{tr}(FO).
\]
Hence all diagonal inequalities are equalities.  If nonnegative
$2\times2$ matrices $A\ge B$ have the same diagonal, then
$\det A\le\det B$.  Applying this observation around the analogous closed
chain of determinants yields equality of all determinants.  For
off-diagonal pairs $(x,y)\ge(x_0,y_0)$, equality of determinants gives
$xy=x_0y_0$; if $x_0y_0>0$, both inequalities must be equalities.
\end{proof}

There is a stronger form in the irreducible branch.

\begin{lemma}[Perron rigidity]\label{lem:perronrigidity}
If $F+T$ is irreducible, then all six carry inequalities are matrix
equalities.  In particular,
\[
 FD=OF,\qquad TO=DF,\qquad TD=DT,\qquad TFD=DF^2.
\tag{K}
\]
\end{lemma}

\begin{proof}
Let $S=Z+O+D$ and let $\Delta_i\ge0$ be the six carry slacks.  Direct
summation gives
\[
 \sum_i\Delta_i=(F+T)S-S(F+T)\ge0.
\]
For positive left and right Perron vectors $x^{\mathsf T},y$ of $F+T$,
the bilinear form of the right side is zero.  A nonnegative matrix between
strictly positive vectors can have zero bilinear form only if it is zero.
Thus every $\Delta_i=0$.  The last identity follows by composition:
$TFD=T(OF)=(TO)F=(DF)F$.
\end{proof}

\subsection{Excluding non-idempotent shapes}

On the complement in Lemma~\ref{lem:complement}, write
\[
 F=\mat{u}{p}{r}{0},\qquad u\ge1,\quad p,r\ge0.
\]
The following two propositions contain the structural core of the proof.

\begin{proposition}[The $u\ge2$ branches]\label{prop:uge2}
No matrices satisfying \eqref{eq:dyncoeff}--\eqref{eq:leftcoeff} and
Lemma~\ref{lem:complement} exist with $u\ge2$.
\end{proposition}

\begin{proof}
There are four zero patterns for $(p,r)$.

If $p=r=0$, Lemma~\ref{lem:rigidity}, $LO\ge LF^2$, and $TG\ge DG$
give $d_{00}=o_{00}\ge u^2$.  Writing the first column of $G$ as
$(c,d)^{\mathsf T}$ with $c>0$, the first component of $TG\ge DG$
forces $d>0$ and $t_{01}>d_{01}$.  If $z_{11}>0$, the lower components of
$FO\ge ZT$ and $TG\ge DG$ reduce the upper-right component of
$TD\ge DT$ to $ud_{01}\ge d_{00}t_{01}$, impossible because
$d_{00}>u$.  If $z_{11}=0$, the carry equalities imply
$u>t_{11}\ge d_{11}$, while the same component gives
\[
 d_{01}(u-t_{11})\ge t_{01}(d_{00}-d_{11}).
\]
Its left side is strictly below $t_{01}(u-t_{11})$ and its right side is
strictly above it.

If $p>0,r=0$, rigidity gives
$z_{10}=o_{10}=0$ and, with $x=o_{00}$,
\[
 x\ge u^2,\qquad ud_{00}+pd_{10}=ux,\qquad
 t_{00}x=ud_{00}.
\]
Thus $t_{00}x+pd_{10}=ux$ and $t_{00}\le u$.  The first component of
$TG\ge DG$ forces the first column $(c,d)^{\mathsf T}$ of $G$ to have
$d>0$ and $t_{01}>d_{01}$.  If $z_{11}>0$, lower carry components force
$T$ and $D$ to have zero lower rows, and the upper-right component of
$TD\ge DT$ becomes $ud_{01}\ge xt_{01}$, impossible since $x>u$.
If $z_{11}=0$, the upper-right component of $FZ\ge ZF$ and the diagonal
carry equations first rule out $t_{00}<u$; hence
$t_{00}=u,d_{10}=0,d_{00}=x,z_{00}=x,t_{10}=0$.  The lower dynamic
component gives $t_{11}\ge d_{11}$.  If $t_{11}=0$, the preceding
impossible inequality returns directly.  Otherwise the summed upper-right
carry equations give $u>t_{11}\ge d_{11}$, whereas $TD\ge DT$ requires
\[
 d_{01}(u-t_{11})\ge t_{01}(x-d_{11}).
\]
The left side is strictly below $t_{01}(u-t_{11})$ and the right side is
strictly above it.

The case $p=0,r>0$ is the transpose-shaped calculation relevant to the first column.  It yields
\[
 t_{01}=0,\quad t_{00}=u,\quad
 z_{01}=o_{01}=d_{01}=0,\quad z_{00}=o_{00}=d_{00},
\]
and then, with $A=au+br>0$,
\[
 u(ao_{00}+bo_{10})\le A d_{00},\qquad
 ao_{00}+bo_{10}\ge uA.
\]
Hence $d_{00}\ge u^2$, whereas the first component of $TG\ge DG$
gives $u\ge d_{00}$.

Finally suppose $p,r>0$.  This is covered by
Proposition~\ref{prop:irreducible}, which in fact permits $u=1$.
The four cases exhaust all nonnegative $(p,r)$.
\end{proof}

\begin{proposition}[Fully off-diagonal branch]\label{prop:irreducible}
There is no solution with $u,p,r\ge1$.
\end{proposition}

\begin{proof}
Now $F+T$ is irreducible, so Lemma~\ref{lem:perronrigidity} applies.  Let
$\lambda>1$ be the Perron root of $F$; the other root $\mu<0$ satisfies
\[
 \lambda+\mu=u,\qquad \lambda\mu=-pr,\qquad
 F^2=uF+prI.
\tag{L}
\]
The matrix $D$ is nonzero because $d_{00}\ge1$.

If $\operatorname{rank}D=1$, write $D=cd^{\mathsf T}$.  The commutation
$TD=DT$ and identity (K), together with Perron--Frobenius for the positive
matrix $F^2$, imply that $d^{\mathsf T}F=\lambda d^{\mathsf T}$ and that
the eigenvalue of $T$ on the image of $D$ is $\lambda$.  The remaining
carry equalities then imply $Fc=\lambda c$.  For the first column $g$ of
$G$, put $h=Dg>0$.  The left and dynamic inequalities yield
\[
 \ell h=\ell Dg\ge\ell FTg\ge\ell FDg
 =\ell Fh=\lambda\ell h,
\]
contradicting $\lambda>1$.

If $\operatorname{rank}D=2$, identities (K) and the carry equalities imply
that $TF$ is similar to $F^2$ and $FT$ is similar to $T^2$.  Traces and
determinants show that $T$ and $F$ share the roots $\lambda,\mu$.
Moreover $FZ=ZF$ gives $Z=\alpha I+\beta F$ with
$\alpha,\beta\ge0$.  Expressing $FTZ=ZT^2$ in an eigenbasis of $F$ shows
that the Perron eigenvector $v>0$ of $F$ is also a Perron eigenvector of
$T$.  Since $TD=DT$, write $Dv=\delta v$.  The dynamic inequality gives
$\delta\le\lambda$, while $LD\ge LFT$ gives
$\delta\ell v\ge\lambda^2\ell v$, so $\delta\ge\lambda^2$.  Again this
contradicts $\lambda>1$.
\end{proof}

\begin{corollary}[Exhaustive residual classification]\label{cor:forms}
Every coefficient assignment not covered by Lemma~\ref{lem:first} has
\[
 F^2=F,\qquad
 F\in\left\{
 \mat{1}{0}{0}{0},
 \mat{1}{p}{0}{0}\ (p\ge1),
 \mat{1}{0}{r}{0}\ (r\ge1)
 \right\}.
\]
\end{corollary}

\begin{proof}
Lemma~\ref{lem:complement} gives $f_{11}=0$ and monotonicity gives
$u=f_{00}\ge1$.  Proposition~\ref{prop:uge2} forces $u=1$, while
Proposition~\ref{prop:irreducible} forces $pr=0$.  Hence
$F=\mat{1}{p}{r}{0}$ with $pr=0$, and direct multiplication gives
$F^2=F$.  The three displayed cases are exhaustive.
\end{proof}

\subsection{Covering the three idempotent forms}

The final step is offset infeasibility.  Each case below is closed by one
of seven parameterized certificate schemes.  Their exact supports and
weights are reconstructed by the artifact; Appendix~\ref{app:certs}
records the mathematical predicates used by the checker.

\begin{proposition}[Lower form]\label{prop:lower}
If $F=\mat{1}{0}{r}{0}$ with $r\ge1$, the relative offset system is
infeasible.
\end{proposition}

\begin{proof}
The lower-triangular calculation in Proposition~\ref{prop:uge2}, up to the
step that used $u\ge2$, remains valid at $u=1$.  It yields
\[
 t_{00}=1,\ t_{01}=0,\quad
 z_{01}=o_{01}=d_{01}=0,\quad z_{00}=o_{00}=d_{00}=1.
\]
The lower-left components of $FD\ge OF$ and $LO\ge LF^2$ force
$o_{10}=r,o_{11}=0$.  The remaining carry and left inequalities force
$z_{11}=t_{11}=0$, $z_{10}=t_{10}=r$, and $d_{10}\ge r$.  These are
exactly the applicability conditions of certificate scheme~5; its strict
right side is $5(a+br)>0$.
\end{proof}

\begin{proposition}[Projection form]\label{prop:projection}
If $F=\operatorname{diag}(1,0)$, one of certificate schemes 2--6 applies.
\end{proposition}

\begin{proof}
Carry rigidity gives $z_{10}=o_{10}=0$ and $d_{00}=o_{00}$.  If
$z_{11}\ge1$, scheme~2 applies.  Otherwise $t_{10}>0$ would force,
successively, $z_{00}=1,z_{01}=0$ and the impossible first-row inequality
$a\ge at_{00}+bt_{10}>a$.  Thus $t_{10}=d_{10}=0$, and the remaining
diagonal equations give $z_{00}=o_{00}=d_{00}=t_{00}=1$.

Set $s=t_{01}$, $y=t_{11}$, $z=z_{01}$, $o=o_{01}$,
$q=o_{11}$, $d=d_{01}$, and $e=d_{11}$.  The carry and left rules imply
\[
 o\ge s+zy,\quad z\ge s+oy,\quad az\ge as+by,\quad qy=0.
\]
If $y>0$, these force $y=1,s=0,o=z$, and $az\ge b$, which is scheme~6.
If $y=0$ and $s\le d$, scheme~5 applies.  If $y=0$ and $s>d$, the last
left inequality gives $ad+be\ge as$, hence $e\ge1$.  For $e=1$ use
scheme~3.  For $e\ge2$ the integer interval of scheme~4 is nonempty; one
valid choice is
\[
 x=\max\left\{\left\lceil b/2\right\rceil,
 \left\lceil a(s-d)/(e+1)\right\rceil\right\}\le b.
\]
These cases are exhaustive.
\end{proof}

\begin{proposition}[Upper form]\label{prop:upper}
If $F=\mat{1}{p}{0}{0}$ with $p\ge1$, one of certificate schemes 4 or 12
applies.
\end{proposition}

\begin{proof}
Lemma~\ref{lem:complement} gives $b>ap$.  The carry rules first imply
$z_{10}=o_{10}=0$.  If $t_{10}>0$, Perron rigidity makes every carry rule
an equality.  Then $z_{11}=0$, $Z=cF$ for $c=z_{00}\ge1$, and, on writing
$T=\mat{x}{y}{s}{w}$ with $s>0$, the carry equations reduce to
\[
 0=(x-1)(x+ps)+ys+p^2s^2,
\]
which is impossible.  Hence $t_{10}=d_{10}=0$, and the diagonal equations
give $z_{00}=o_{00}=d_{00}=t_{00}=1$.

Put $s=t_{01}$, $w=t_{11}$, $z=z_{01}$, $k=z_{11}$,
$o=o_{01}$, $q=o_{11}$, $d=d_{01}$, and $e=d_{11}$.  The required
scalar consequences are
\begin{gather*}
 z+pk\ge p,\quad o+pq\ge s+zw,\quad d+pe\ge p,\\
 z+sk\ge s+ow,\quad o+sq\ge p,\quad d+se\ge s+dw,\\
 az+bk\ge as+bw,\quad ao+bq\ge ap,\quad
 ad+be\ge a(s+pw),\quad kw=qw=0.
\tag{N}
\end{gather*}
If $w>0$, (N) forces $k=q=0$, then $w=1,s=0,z=o$.  If $e\le1$,
scheme~12 applies; if $e\ge2$, scheme~4 applies with $x=b$.  If $w=0$,
scheme~12 again applies for $e\le1$, while for $e\ge2$ scheme~4 applies
with $x=b$, using $ap-b<0$.  Thus every upper-form assignment is covered.
\end{proof}

\begin{theorem}[Unbounded two-dimensional no-go]\label{thm:main}
The system $\Tsys$ has no relative natural affine matrix interpretation of
dimension two in which the two dynamic rules are strict and the other nine
rules are weak.  Matrix entries and offsets are unrestricted elements of
$\N$.
\end{theorem}

\begin{proof}
If a coefficient inequality
\eqref{eq:dyncoeff}--\eqref{eq:leftcoeff} fails, no interpretation exists.
Otherwise, Lemma~\ref{lem:first} supplies a certificate when
$\ell F\ge\ell$.  On its complement,
Corollary~\ref{cor:forms} gives exactly the lower, projection, or upper
idempotent form.  Propositions~\ref{prop:lower}--\ref{prop:upper} cover
those forms by exact nonnegative Farkas certificates with positive strict
right side.  Lemma~\ref{lem:farkas} excludes the offsets in every case.
\end{proof}

\begin{corollary}
No such interpretation exists in dimension at most two.
\end{corollary}

\section{Automation and reproducibility}

The implementation has two deliberately separated layers.

\paragraph{Discovery.}
Bounded enumeration and linear programming were used to find small supports
for dual certificates.  Counterexamples to provisional families enlarged
the search region.  The numerical LP output is not used by the final
verifier.

\paragraph{Exact checking.}
The Python module \texttt{sag\_collatz.matrix\_attack} stores the original
rewrite rules, constructs word matrices and offsets, and validates all
certificate rows with arbitrary-precision integer arithmetic.  The only
certificate with a free integer weight solves its fourteen inequalities as
an exact integer interval rather than by bounded trial.  A separate C++17
program exhaustively regression-tests bounded coefficient boxes.

\paragraph{What the checker proves.}
For any supplied matrix assignment, the checker validates the coefficient
obligations, selects a parameterized family, and reconstructs a Farkas
contradiction.  It does not by itself prove that the families cover every
unbounded assignment.  That universal coverage is the content of
Lemmas~\ref{lem:complement}--\ref{lem:perronrigidity},
Propositions~\ref{prop:uge2}--\ref{prop:upper}, and
Corollary~\ref{cor:forms}.  We call the result \emph{computer-assisted}, not
formally verified in Lean, Coq, or an SMT proof calculus.

The repository contains unit tests for rule expansion, exact certificate
validation, adversarial mutations, the scalar theorem, bounded exhaustive
counts, and structural families at coefficients far outside the discovery
boxes.  Reproduction commands and expected outputs are recorded in
\texttt{ARTIFACT.md}.

\section{Related work}

Matrix interpretations for term and string rewriting were developed as a
general automatic termination technique by Endrullis, Waldmann, and
Zantema~\cite{endmat}, building on automatic string-rewriting termination
methods such as~\cite{zantema}.  Relative termination and rule removal are
standard ways to combine weak and strict orientation obligations
\cite{termbook}.

YAH~\cite{yah} provide the direct point of departure: they prove the
equivalence between Collatz and termination of $\Tsys$, automatically find
interpretations for weakened systems, and identify nonexistence of suitable
interpretations for the full system as a useful possibility to investigate.
They also prove a substantially broader natural-matrix obstruction for a
different, previously studied \emph{unary} Collatz rewrite encoding, even in
combination with dependency pairs.  That theorem does not apply to the
mixed-base system $\Tsys$ considered here; conversely, our theorem is limited
to dimensions at most two but addresses the mixed-base encoding that makes
their automated positive results possible.
Our work is complementary to automatic synthesis.  It certifies that one
specific synthesis space---one-stage natural affine interpretations in
dimensions at most two---is empty even with unbounded coefficients.

The result is not a lower bound for all matrix interpretations.  Natural
matrices in arbitrary dimension are already only one termination language;
arctic and tropical semirings, dependency pairs, and richer interpretation
schemes have different expressive power.  Nor does failure in every fixed
finite dimension automatically imply nontermination of the rewrite system.

\section{Limitations and conclusion}

Theorem~\ref{thm:main} is a no-go theorem about a proof method.  It neither
proves nor refutes Collatz.  Specifically, it leaves open:
\begin{itemize}[leftmargin=*]
  \item natural matrix interpretations of dimension at least three;
  \item arctic, tropical, and other ordered semirings;
  \item dependency-pair transformations;
  \item multi-stage rule-removal sequences outside the single relative
        obligation studied here;
  \item alternative rewrite encodings equivalent to Collatz.
\end{itemize}

Within its scope, however, the theorem is complete: there is no coefficient
ceiling, no offset ceiling, and no unclassified zero pattern.  The proof also
demonstrates a useful workflow for automated reasoning.  Small exhaustive
search can expose recurring dual supports; exact arithmetic can turn them
into trusted local certificates; and structural algebra can remove the
finite search boundary.  Here that process closes a concrete low-dimensional
search space proposed by an influential automated approach to Collatz.

\section*{Statements and Declarations}

\paragraph{Competing interests.}
The authors declare no competing interests.

\paragraph{Funding.}
No funding was received for this work.

\paragraph{Data and code availability.}
All source code, exact certificates, tests, and the manuscript sources are
publicly available at
\url{https://github.com/kadirbekov-ship-it/collatz-matrix-no-go} and in a
versioned archival deposit at
\url{https://doi.org/10.5281/zenodo.22098492}.

\paragraph{Author contributions.}
K.K. and D.K. contributed to the conception, oversight of the AI-assisted
mathematical analysis, verification workflow, and review of the manuscript.
Both authors reviewed and approved the submitted version.

\paragraph{Use of generative AI.}
OpenAI Codex using GPT-5.6 Sol was used as an adversarial mathematical and
coding assistant for exploratory derivations, implementation, testing, and
language drafting.  Anthropic Claude Opus was used for a separate adversarial
audit of intermediate claims and calculations.  All theorem statements,
proofs, computations, citations, and the final manuscript were reviewed and
accepted by the human authors.  Model outputs were not treated as independent
evidence, and the systems are not authors.

\appendix
\section{Certificate predicates}\label{app:certs}

For auditability we summarize the seven schemes used by the universal proof.
Each predicate is sufficient for the exact row emitted by the implementation;
the checker expands the row against all fourteen offset variables.

Write $\Delta_r^i$ for the offset difference of rule $r$ in component $i$.
The seven nonnegative combinations are
\begin{align*}
C_1={}&a\Delta_{\mathrm{odd}}^0+b\Delta_{\mathrm{odd}}^1
       +\Delta_{\mathrm{left2}}^0,\\
C_2={}&a\Delta_{\mathrm{even}}^0+a\Delta_{\mathrm{odd}}^0
       +b\Delta_{\mathrm{odd}}^1+b\Delta_{\mathrm{f1}}^1
       +\Delta_{\mathrm{left2}}^0,\\
C_3={}&a\Delta_{\mathrm{even}}^0+b\Delta_{\mathrm{even}}^1
       +a\Delta_{\mathrm{odd}}^0+b\Delta_{\mathrm{t2}}^1
       +\Delta_{\mathrm{left2}}^0,\\
C_4(x)={}&a\Delta_{\mathrm{even}}^0+x\Delta_{\mathrm{even}}^1
       +a\Delta_{\mathrm{odd}}^0+x\Delta_{\mathrm{odd}}^1
       +x\Delta_{\mathrm{t2}}^1+\Delta_{\mathrm{left2}}^0.
\end{align*}
For scheme~5 put $A=a+bf_{10}$; then
\begin{align*}
C_5={}&2A\Delta_{\mathrm{even}}^0+3A\Delta_{\mathrm{odd}}^0
 +a\Delta_{\mathrm{f1}}^0+b\Delta_{\mathrm{f1}}^1\\
& +(2a+bf_{10})\Delta_{\mathrm{f2}}^0
 +b\Delta_{\mathrm{f2}}^1+b\Delta_{\mathrm{t2}}^1
 +\Delta_{\mathrm{left0}}^0+\Delta_{\mathrm{left1}}^0
 +\Delta_{\mathrm{left2}}^0.
\end{align*}
Scheme~6 is
\begin{align*}
C_6={}&3a\Delta_{\mathrm{even}}^0+3a\Delta_{\mathrm{odd}}^0
 +a\Delta_{\mathrm{f1}}^0+a\Delta_{\mathrm{f2}}^0
 +3b\Delta_{\mathrm{f2}}^1\\
&+a\Delta_{\mathrm{t0}}^0+2b\Delta_{\mathrm{t1}}^1
 +\Delta_{\mathrm{left1}}^0+2\Delta_{\mathrm{left2}}^0.
\end{align*}
For scheme~12 put
\[
 K=az_{01}+bz_{11},\qquad
 q=a\bigl(o_{01}-f_{01}(1-o_{11})\bigr),\qquad
 H=af_{01}+d_{11}K+b.
\]
Then
\begin{align*}
C_{12}={}&q\Delta_{\mathrm{even}}^1+3a\Delta_{\mathrm{odd}}^0
 +H\Delta_{\mathrm{odd}}^1+a\Delta_{\mathrm{f1}}^0
 +b\Delta_{\mathrm{f1}}^1\\
&+a\Delta_{\mathrm{f2}}^0+af_{01}\Delta_{\mathrm{f2}}^1
 +\Delta_{\mathrm{left0}}^0+\Delta_{\mathrm{left2}}^0.
\end{align*}
In every branch used in the proof, the displayed weights are nonnegative,
all fourteen coefficients of the relevant $C_i$ are nonpositive, and the
strict right side is positive.  These assertions follow by direct expansion
from the branch inequalities; the implementation performs the same expansion
without simplification.

Scheme~1 is Lemma~\ref{lem:first}.  In the complement
$b>0,f_{11}=0,af_{01}<b$, scheme~2 is equivalent to
\[
 f_{10}=0,\qquad z_{11}\ge1,\qquad
 ad_{01}+bd_{11}\ge a(f_{01}+t_{01}).
\]
Scheme~5 is equivalent there to
\begin{gather*}
 f_{00}=1,\quad f_{01}=0,\quad t_{11}=0,\quad t_{10}\le f_{10},\\
 2a(1-o_{00})+bf_{10}(2-o_{00})-bo_{10}\le0,\\
 a(1-z_{00})+b(2f_{10}-d_{10}-z_{10})\le0,\\
 t_{00}\le d_{00},\qquad t_{01}\le d_{01}.
\end{gather*}
Schemes~3 and~6 are the special projection-form branches identified in
Proposition~\ref{prop:projection}; scheme~12 is the upper-form branch in
Proposition~\ref{prop:upper}.

Scheme~4 contains one positive integer parameter $x$.  Every combined
coefficient has the form $A_j+xB_j$.  The implementation intersects the
exact lower and upper integer bounds
\[
 x\ge\max_{B_j<0}\left\lceil\frac{A_j}{-B_j}\right\rceil,
 \qquad
 x\le\min_{B_j>0}\left\lfloor\frac{-A_j}{B_j}\right\rfloor,
\]
together with $x\ge1$, and then reconstructs the certificate.  Thus no
finite search limit is hidden in the parameterized family.

The machine-readable supports and all unsimplified coefficient predicates
are distributed with the artifact as an additional check against transcription
errors.

\begin{thebibliography}{9}

\bibitem{yah}
E. Yolcu, S. Aaronson, M. J. H. Heule,
An automated approach to the Collatz conjecture,
\emph{Journal of Automated Reasoning} 67, 15 (2023).
\url{https://doi.org/10.1007/s10817-022-09658-8}

\bibitem{endmat}
J. Endrullis, J. Waldmann, H. Zantema,
Matrix interpretations for proving termination of term rewriting,
\emph{Journal of Automated Reasoning} 40, 195--220 (2008).
\url{https://doi.org/10.1007/s10817-007-9087-9}

\bibitem{zantema}
H. Zantema,
Termination of string rewriting proved automatically,
\emph{Journal of Automated Reasoning} 34, 105--139 (2005).
\url{https://doi.org/10.1007/s10817-005-6545-0}

\bibitem{schrijver}
A. Schrijver,
\emph{Theory of Linear and Integer Programming},
Wiley, Chichester (1986).

\bibitem{termbook}
F. Baader, T. Nipkow,
\emph{Term Rewriting and All That},
Cambridge University Press, Cambridge (1998).
\url{https://doi.org/10.1017/CBO9781139172752}

\end{thebibliography}

\end{document}
